#I. Relational Operators x DataCamp
# ==,!=, &, |, !(e.g is.numeric(), !is.numeric("hey"))
linked_in <- c(16, 9,16)
linked_in > 10 #returns TRUE FALSE TRUE
"goodbye" != "hello" #returns false

#& vs && and | vs ||
#Illustration:
#c(TRUE, TRUE, FALSE) & c(TRUE, FALSE, FALSE) #returns TRUE FALSE FALSE but
#c(TRUE, TRUE, FALSE) & c(TRUE, FALSE, FALSE) #returns only TRUE, only access the first element
#Accessing data frame 2nd column example: name <- li_df[,2]


#II. Conditional Statements in R x DataCamp





#III. While Loop in R x Datacamp
while(condition){
  expr  
} #somewhat same as C

while(ctr <= 7){
  print(paste(""counter is set to ", ctr)) #just like in C, paste is used to concatenate stings and numbers yada yada
  ctr  <- ctr + 1 #how much does ctr will increment per loop
}
if we just set ctr <- 1 before loop and we didn't use the increment concept (ctr <- ctr+1), then the program will run indefinitely

#break statement:
ctr <- 1
while(ctr <=7){
   if (ctr%%5==0){ #it breaks wen R finds it 
   break
   }
   #otherwise 
   print(paste("ctr is set to", ctr))
   ctr <- ctr + 1
}

#e.g.
# Initialize the speed variable
speed <- 64

# Extend/adapt the while loop
while (speed > 30) {
  print(paste("Your speed is", speed))
  if (speed > 48) {
    print("Slow down big time!")
    speed <- speed - 11
  } else {
    print("Slow down!")
    speed <- speed - 6
  }
}
#Output: 
[1] "Your speed is 64"
[1] "Slow down big time!"
[1] "Your speed is 53"
[1] "Slow down big time!"
[1] "Your speed is 42"
[1] "Slow down!"
[1] "Your speed is 36"
[1] "Slow down!"


#FOR LOOP x DATA_CAMP
e.g. new_york <- c("newyork", "Thailand", "Luzon")
for(cities in new_york){
  print(cities)
} #also works in lists.

#break statement - a statement that breaks the for loop altogether
e.g  new_york <- c("newyork", "Thailand", "Luzon")
      for(cities in new_york){
          if(nchar(cities) == 6){
            break
          }
          print(cities)
      }
#next statement - skips the current loop and proceed to another
e.g  new_york <- c("newyork", "Thailand", "Luzon")
      for(cities in new_york){
          if(nchar(cities) == 6){
            next 
          }
          print(cities)
      }





